\section{September 22, 2026}

Throughout, $(X,\mathcal M,\mu)$ is a measure space, and
\[
  L^+=\{f\colon X\to[0,+\infty]:f\text{ is measurable}\}.
\]
We continue to use $\mathbf 1_E$ for the characteristic function of a
measurable set $E$ and the convention $0\cdot\infty=0$.

\subsection{Nonnegative integrals and infinite values}

Recall that, for a nonnegative simple function in standard form,
\[
  \phi=\sum_{j=1}^N a_j\mathbf 1_{E_j},
  \qquad
  \int_X\phi\,d\mu=\sum_{j=1}^N a_j\mu(E_j).
\]
Moreover, $A\mapsto\int_A\phi\,d\mu$ is a measure. For $f\in L^+$,
\[
  \int_X f\,d\mu
  =\sup\left\{\int_X\phi\,d\mu:
    \phi\text{ is simple},\ 0\leq\phi\leq f\right\}.
\]
In particular, $f\leq g$ implies $\int_X f\,d\mu\leq\int_X g\,d\mu$:
every simple function admissible for $f$ is also admissible for $g$.
For $c\geq0$, the same definition gives
$\int_X cf\,d\mu=c\int_X f\,d\mu$. For $c>0$, use the correspondence
$\phi\mapsto c\phi$ between the two classes of simple functions; the
case $c=0$ follows from our convention.

\begin{example}[Infinite values and integrals]
  Let $m$ be Lebesgue measure on $\bb R$, and let $C$ be the
  middle-thirds Cantor set.
  \begin{enumerate}[label=(\roman*)]
    \item The function $f=+\infty\,\mathbf 1_C$ has integral zero.
      Indeed, if $0\leq\phi\leq f$ is simple, all its positive level
      sets are contained in the null set $C$, so $\int\phi\,dm=0$.
    \item For $f=+\infty\,\mathbf 1_{[0,1]}$, the simple function
      $N\mathbf 1_{[0,1]}$ lies below $f$ for every $N\in\bb N$.
      Consequently $\int f\,dm\geq N$ for every $N$, so
      $\int f\,dm=+\infty$.
    \item Define $f(x)=x^{-2}$ for $0<|x|\leq1$, set $f(0)=+\infty$,
      and set $f(x)=0$ for $|x|>1$. For $0<a<1$,
      \[
        f\geq a^{-2}\mathbf 1_{(0,a)},
        \qquad \int_{\bb R}f\,dm\geq a^{-1}.
      \]
      Thus $\int f\,dm=+\infty$, even though $f$ is finite almost
      everywhere. The divergence comes from its values near zero,
      not merely from its value at zero.
    \item With the same $f$ but with Dirac measure $\delta_0$,
      \[
        \int_{\bb R}f\,d\delta_0=f(0)=+\infty.
      \]
      Whether a single point contributes to an integral depends on
      the measure.
  \end{enumerate}
\end{example}

\subsection{The monotone convergence theorem}

\begin{theorem}[Monotone convergence theorem]
  \label{thm:mct-sep-22}
  Suppose $f_n\in L^+$ and $f_n\uparrow f$ pointwise. Then $f\in L^+$
  and
  \[
    \int_X f\,d\mu=\lim_{n\to\infty}\int_X f_n\,d\mu.
  \]
  Both sides are allowed to be $+\infty$.
\end{theorem}

\begin{proof}
  Measurability of $f$ follows from
  Proposition~\ref{prop:measurable-sequential-operations-sep-17}.
  By monotonicity, the sequence of integrals is increasing and
  \[
    A:=\lim_{n\to\infty}\int_X f_n\,d\mu\leq\int_X f\,d\mu.
  \]
  To prove the reverse inequality, fix a simple function
  $0\leq\phi\leq f$ and a number $0<\alpha<1$. Set
  \[
    E_n=\{x\in X:f_n(x)\geq\alpha\phi(x)\}.
  \]
  These sets are measurable: on each level set $\{\phi=a_j\}$, the
  defining condition is the measurable inequality
  $f_n\geq\alpha a_j$. Since $f_n$ increases, $E_n\subseteq E_{n+1}$.
  In fact $E_n\uparrow X$. If $\phi(x)=0$, then $x\in E_n$ for every
  $n$. If $\phi(x)>0$, then
  \[
    \lim_n f_n(x)=f(x)\geq\phi(x)>\alpha\phi(x),
  \]
  so $x\in E_n$ for all sufficiently large $n$.

  Now $\alpha\phi\mathbf 1_{E_n}\leq f_n$, and therefore
  \[
    \int_X f_n\,d\mu\geq\alpha\int_{E_n}\phi\,d\mu.
  \]
  Apply continuity from below to the measure
  $\nu(E)=\int_E\phi\,d\mu$, established in
  Proposition~\ref{prop:simple-integral-properties-sep-17}. This gives
  \[
    A\geq\alpha\lim_n\nu(E_n)
      =\alpha\nu(X)=\alpha\int_X\phi\,d\mu.
  \]
  If $\int\phi\,d\mu=+\infty$, this already forces $A=+\infty$.
  Otherwise, let $\alpha\uparrow1$ to obtain $A\geq\int\phi\,d\mu$.
  Taking the supremum over all such $\phi$ proves
  $A\geq\int f\,d\mu$.
\end{proof}

\begin{remark}
  The factor $\alpha<1$ is essential in this proof: $f_n(x)$ can
  approach $\phi(x)=f(x)$ from below without ever reaching it. It
  creates the strict gap needed to ensure $\bigcup_n E_n=X$.
\end{remark}

\begin{figure}[htbp]
  \centering
  \begin{tikzpicture}[x=1.55cm,y=1.05cm,>=Latex]
    \fill[blue!8] (0,0) -- plot[domain=0:5,samples=60]
      (\x,{0.55+0.8*sin(36*\x)}) -- (5,0) -- cycle;
    \draw[->] (-0.1,0)--(5.5,0) node[right] {$x$};
    \draw[->] (0,-0.1)--(0,3.25) node[above] {value};
    \draw[blue!45!black,thick,domain=0:5,samples=60]
      plot (\x,{0.55+0.8*sin(36*\x)});
    \draw[blue!65!black,thick,domain=0:5,samples=60]
      plot (\x,{1.1+0.8*sin(36*\x)});
    \draw[blue!85!black,thick,domain=0:5,samples=60]
      plot (\x,{1.65+0.8*sin(36*\x)});
    \draw[black,very thick,domain=0:5,samples=60]
      plot (\x,{2.2+0.8*sin(36*\x)});
    \node[right] at (5,0.55) {$f_1$};
    \node[right] at (5,1.1) {$f_2$};
    \node[right] at (5,1.65) {$f_3$};
    \node[right] at (5,2.2) {$f$};
    \draw[->,red!70!black,thick] (2.5,1.5)--(2.5,2.92);
  \end{tikzpicture}
  \caption{For an increasing sequence of nonnegative functions, the
    regions below the graphs increase to the region below the limit.
    The monotone convergence theorem makes this area intuition valid
    on an arbitrary measure space.}
  \label{fig:mct-sep-22}
\end{figure}

\begin{corollary}[Computing with simple approximations]
  If $0\leq\phi_n\uparrow f$ are simple, then
  \[
    \int_X f\,d\mu=\lim_n\int_X\phi_n\,d\mu.
  \]
  In particular, the approximations constructed in
  Theorem~\ref*{thm:simple-approximation-sep-17} compute the integral.
\end{corollary}

\begin{proposition}[Integrating sums of nonnegative functions]
  \label{prop:nonnegative-series-sep-22}
  If $f_n\in L^+$, then
  \[
    \int_X\left(\sum_{n=1}^{\infty}f_n\right)\,d\mu
      =\sum_{n=1}^{\infty}\int_X f_n\,d\mu.
  \]
  The analogous identity holds for finite sums.
\end{proposition}

\begin{proof}
  First take two functions $f,g\in L^+$. Choose nonnegative simple
  functions $\phi_n\uparrow f$ and $\psi_n\uparrow g$. Then
  $\phi_n+\psi_n\uparrow f+g$, so the monotone convergence theorem and
  additivity of the simple integral give
  \begin{align*}
    \int_X(f+g)\,d\mu
      &=\lim_n\int_X(\phi_n+\psi_n)\,d\mu\\
      &=\lim_n\left(\int_X\phi_n\,d\mu+\int_X\psi_n\,d\mu\right)\\
      &=\int_X f\,d\mu+\int_X g\,d\mu.
  \end{align*}
  All quantities are nonnegative, so no undefined subtraction occurs
  even when a limit is infinite. Induction proves finite additivity.
  Finally, apply monotone convergence to
  $S_N=\sum_{n=1}^N f_n\uparrow\sum_{n=1}^{\infty}f_n$:
  \[
    \int_X\sum_{n=1}^{\infty}f_n\,d\mu
      =\lim_N\sum_{n=1}^N\int_X f_n\,d\mu
      =\sum_{n=1}^{\infty}\int_X f_n\,d\mu.
  \]
\end{proof}

\subsection{Null sets and almost everywhere convergence}

\begin{lemma}[Zero integral]
  \label{lem:zero-integral-sep-22}
  For $f\in L^+$,
  \[
    \int_X f\,d\mu=0
    \quad\Longleftrightarrow\quad f=0\quad\mu\text{-a.e.}
  \]
\end{lemma}

\begin{proof}
  If $f=0$ a.e., every nonnegative simple function $\phi\leq f$
  vanishes a.e. Its positive level sets are null, so
  $\int\phi\,d\mu=0$. Taking the supremum gives $\int f\,d\mu=0$.

  Conversely, suppose $\int f\,d\mu=0$. For each $n\geq1$, set
  $A_n=\{f>1/n\}$. Since $f\geq n^{-1}\mathbf 1_{A_n}$,
  \[
    0=\int_X f\,d\mu\geq\frac1n\mu(A_n),
  \]
  whence $\mu(A_n)=0$. The identity
  $\{f>0\}=\bigcup_{n=1}^{\infty}A_n$ now proves the claim.
\end{proof}

\begin{corollary}[Changing a function on a null set]
  \label{cor:ae-integrals-sep-22}
  If $f,g\in L^+$ and $f=g$ a.e., then $\int f\,d\mu=\int g\,d\mu$.
  More generally, $f\leq g$ a.e. implies
  $\int f\,d\mu\leq\int g\,d\mu$.
\end{corollary}

\begin{proof}
  Choose a measurable null set $N$ outside which the equality or
  inequality holds. The lemma gives
  $\int f\mathbf 1_N\,d\mu=\int g\mathbf 1_N\,d\mu=0$.
  Split each function into its parts on $N$ and $X\setminus N$, use
  additivity, and compare the parts on $X\setminus N$.
\end{proof}

\begin{corollary}[Almost everywhere monotone convergence]
  If $f_n,f\in L^+$ and $f_n\uparrow f$ a.e., then
  \[
    \lim_n\int_X f_n\,d\mu=\int_X f\,d\mu.
  \]
\end{corollary}

\begin{proof}
  Outside one measurable null set $N$, both monotonicity and convergence
  hold. If these properties are initially stated separately a.e., take
  the union of their countably many exceptional null sets. Then
  $f_n\mathbf 1_{X\setminus N}\uparrow f\mathbf 1_{X\setminus N}$
  everywhere. Apply Theorem~\ref{thm:mct-sep-22} and
  Corollary~\ref{cor:ae-integrals-sep-22}.
\end{proof}

\subsection{Fatou's lemma}

Pointwise convergence by itself does not justify interchanging a limit
and an integral.

\begin{example}[Two ways for mass to disappear]
  \label{ex:escaping-mass-sep-22}
  For Lebesgue measure on $\bb R$, consider
  \[
    u_n=\mathbf 1_{[n,n+1]},
    \qquad v_n=n\mathbf 1_{(0,1/n)}.
  \]
  Both sequences converge to zero at every point, but for every $n$,
  \[
    \int u_n\,dm=\int v_n\,dm=1.
  \]
  In the first example
  the support moves to infinity; in the second it shrinks toward zero
  while the height increases. Neither sequence is increasing.
\end{example}

\begin{figure}[htbp]
  \centering
  \begin{tikzpicture}[>=Latex,font=\small]
    \begin{scope}[x=0.65cm,y=1.1cm]
      \draw[->] (0,0)--(6.7,0) node[right] {$x$};
      \draw[->] (0,0)--(0,1.65);
      \foreach \a/\col in {1/blue,3/teal,5/orange} {
        \fill[\col!15] (\a,0) rectangle (\a+1,1);
        \draw[\col!75!black,thick]
          (\a,0)--(\a,1)--(\a+1,1)--(\a+1,0);
      }
      \node[left] at (0,1) {$1$};
      \node[below] at (3,0) {$n$};
      \draw[->,thick] (1.5,1.35)--(5.5,1.35);
      \node at (3.3,-0.8) {support moves to infinity};
    \end{scope}
    \begin{scope}[shift={(6.1cm,0)},x=2.6cm,y=0.43cm]
      \draw[->] (0,0)--(1.35,0) node[right] {$x$};
      \draw[->] (0,0)--(0,4.6);
      \foreach \n/\col in {1/blue,2/teal,4/orange} {
        \draw[\col!75!black,thick]
          (0,\n)--(1/\n,\n)--(1/\n,0);
      }
      \node[left] at (0,4) {$n$};
      \node[below] at (0.25,0) {$1/n$};
      \node at (0.6,-2.05) {height grows as support shrinks};
    \end{scope}
  \end{tikzpicture}
  \caption{Each rectangle has area $1$, but both pointwise limits are
    zero. Endpoint values do not affect these Lebesgue integrals.}
  \label{fig:escaping-mass-sep-22}
\end{figure}

\begin{theorem}[Fatou's lemma]\label{thm:fatou-sep-22}
  For any sequence $f_n\in L^+$,
  \[
    \int_X\liminf_{n\to\infty}f_n\,d\mu
      \leq\liminf_{n\to\infty}\int_X f_n\,d\mu.
  \]
  In particular, if $f_n\to f$ a.e. and $f\in L^+$, then
  \[
    \int_X f\,d\mu\leq\liminf_{n\to\infty}\int_X f_n\,d\mu.
  \]
\end{theorem}

\begin{proof}
  Let $g_k=\inf_{n\geq k}f_n$. These functions are measurable,
  nonnegative, and increase to $\liminf_n f_n$. Hence
  \[
    \int_X\liminf_n f_n\,d\mu
      =\lim_k\int_X g_k\,d\mu
  \]
  by monotone convergence. For every $n\geq k$, $g_k\leq f_n$, so
  \[
    \int_X g_k\,d\mu\leq\inf_{n\geq k}\int_X f_n\,d\mu.
  \]
  Letting $k\to\infty$ proves the inequality. In the almost everywhere
  convergence case, $\liminf_n f_n=f$ a.e.; changing a function on a
  null set does not change its integral.
\end{proof}

\begin{example}[Fatou's lemma for series]
  On $\bb N$ with counting measure, the nonnegative integral is a sum:
  $\int h\,d\mu=\sum_{k=1}^{\infty}h(k)$. This follows by applying
  monotone convergence to $h\mathbf 1_{\{1,\ldots,N\}}$ (with
  truncation as well if $h$ takes infinite values). Thus if
  $a_{n,k}\geq0$ and $a_{n,k}\to a_k$ for each fixed $k$, then
  \[
    \sum_{k=1}^{\infty}a_k
      \leq\liminf_{n\to\infty}\sum_{k=1}^{\infty}a_{n,k}.
  \]
\end{example}

\FloatBarrier

\subsection{Integrating real-valued functions}

Recall $f^+=\max\{f,0\}$ and $f^-=\max\{-f,0\}$, so that
$f=f^+-f^-$ and $|f|=f^++f^-$. Both parts are nonnegative measurable
functions whenever $f\colon X\to\bb R$ is measurable.

\begin{definition}[Signed integral and integrability]
  If at least one of $\int f^+\,d\mu$ and $\int f^-\,d\mu$ is finite,
  define
  \[
    \int_X f\,d\mu=\int_X f^+\,d\mu-\int_X f^-\,d\mu.
  \]
  If both are finite, $f$ is called \emph{integrable}. Equivalently,
  \[
    f\text{ is integrable}\quad\Longleftrightarrow\quad
    \int_X|f|\,d\mu<\infty,
  \]
  since $\int|f|=\int f^++\int f^-$. If both parts have infinite
  integral, $\int f$ is undefined: the expression
  $+\infty-(+\infty)$ has no assigned value.
\end{definition}

Write $\mathcal L^1(X,\mu;\bb R)$ for the real-valued integrable
measurable functions. In the next lecture, $L^1$ will denote the space
obtained by identifying functions that agree almost everywhere.

\begin{proposition}[Linearity of the real integral]
  \label{prop:real-linearity-sep-22}
  The set $\mathcal L^1(X,\mu;\bb R)$ is a real vector space, and
  integration is a linear functional on it. Thus, for integrable
  $f,g$ and $a,b\in\bb R$,
  \[
    \int_X(af+bg)\,d\mu=a\int_X f\,d\mu+b\int_X g\,d\mu.
  \]
\end{proposition}

\begin{proof}
  First, $|af+bg|\leq|a||f|+|b||g|$, so $af+bg$ is integrable.
  If $c\geq0$, then $(cf)^+=cf^+$ and $(cf)^-=cf^-$, giving
  $\int cf=c\int f$. If $c<0$, the positive and negative parts switch:
  $(cf)^+=(-c)f^-$ and $(cf)^-=(-c)f^+$, giving the same identity.

  For additivity, put $h=f+g$. The pointwise identity
  \[
    h^++f^-+g^-=h^-+f^++g^+
  \]
  involves only nonnegative functions. Integrating and using
  nonnegative additivity gives
  \[
    \int h^++\int f^-+\int g^-
      =\int h^-+\int f^++\int g^+.
  \]
  All six integrals are finite, so rearrangement is legitimate and
  yields $\int h=\int f+\int g$.
\end{proof}

\begin{remark}
  If $f,g$ are real-valued integrable functions and $f\leq g$ a.e.,
  then $\int f\,d\mu\leq\int g\,d\mu$. Indeed, $g-f$ is
  nonnegative a.e., so its negative part has integral zero, and
  linearity gives $\int g-\int f=\int(g-f)\geq0$.
\end{remark}

\FloatBarrier
